
\documentclass[11pt]{article}
\usepackage{amsmath,amssymb,amsthm}
\newtheorem{conjecture}{Conjecture}
\newtheorem{theorem}{Theorem}

\title{Two Parametric Families of Polynomial Continued Fractions\\
for Reciprocals of Partial Logarithms and Multiples of $1/\pi$}
\author{Discovery via Ramanujan Breakthrough Generator}
\date{April 2026}

\begin{document}
\maketitle

\begin{abstract}
We present two infinite parametric families of polynomial continued fractions (PCFs)
discovered through systematic numerical search using a modular signature methodology:
a \emph{logarithmic ladder} producing $1/\ln(k/(k-1))$ for all real $k > 1$,
and a \emph{pi family} producing $2^{2m+1}/(\pi \binom{2m}{m})$ for all integers $m \geq 0$.
Both families are verified to 500+ matching decimal digits via independent
high-precision evaluation and PSLQ integer relation detection, with select cases
\textbf{certified to 1200+ digits via Arb ball arithmetic}.
The logarithmic family admits a natural hypergeometric explanation via
$\ln(1+x) = x \cdot {}_2F_1(1,1;2;-x)$, while extending to all real parameters $k > 1$
(verified for 15 values including non-integers to $90$--$120$ digits).
The pi family exhibits a striking parity phenomenon: odd parameters yield
transcendental multiples of $1/\pi$ involving central binomial coefficients,
while even parameters collapse to Wallis-type rationals.
\end{abstract}

\section{Main Results}

\begin{conjecture}[Logarithmic Ladder]\label{conj:log}
For all integers $k \geq 2$, the polynomial continued fraction
\[
  k + \cfrac{-k \cdot 1^2}{(k+2) + \cfrac{-k \cdot 2^2}{(k+4) + \cfrac{-k \cdot 3^2}{(k+6) + \cdots}}}
\]
with $a(n) = -kn^2$ and $b(n) = (k+1)n + k$ converges to
\[
  \frac{1}{\ln\!\left(\frac{k}{k-1}\right)}.
\]
\end{conjecture}

\noindent Numerically verified for integer $k = 2, 3, \ldots, 9$ to $118$--$120$
matching decimal digits and for 15 real values of $k$ (including $k = 1.5, 2.5, 3.5, \ldots$)
to $90$--$120$ digits. Cases $k = 2$ and $k = 3$ certified via Arb ball
arithmetic to $\mathbf{1208}$ and $\mathbf{1817}$ digits respectively.

\smallskip
\noindent\textbf{Hypergeometric connection.} Using the identity
$\ln(1+x) = x \cdot {}_2F_1(1,1;2;-x)$ with $x = 1/(k-1)$,
the logarithmic ladder is the reciprocal of a scaled Gauss hypergeometric evaluation.
This provides a natural framework for a potential proof via the Gauss continued
fraction for ${}_2F_1(1,1;2;z)$.

\smallskip

\begin{conjecture}[Pi Family]\label{conj:pi}
For all non-negative integers $m$, the PCF with
\[
  a(n) = -n(2n - (2m+1)), \quad b(n) = 3n + 1
\]
converges to
\[
  \frac{2^{2m+1}}{\pi \binom{2m}{m}}.
\]
\end{conjecture}

\noindent Verified for $m = 0, 1, \ldots, 7$ to $99$--$119$ matching digits.
Cases $m = 0$ ($2/\pi$) and $m = 1$ ($4/\pi$) certified via Arb to
$\mathbf{1207}$ and $\mathbf{1217}$ digits respectively.
The closed form can be rewritten using the Gamma function as
$\text{val}(m) = 2\,\Gamma(m+1) / (\sqrt{\pi}\,\Gamma(m+\tfrac{1}{2}))$.
The first two cases are:
\begin{align*}
  m = 0: &\quad a(n) = -n(2n-1),\ b(n) = 3n+1 \longrightarrow \frac{2}{\pi}, \\
  m = 1: &\quad a(n) = -n(2n-3),\ b(n) = 3n+1 \longrightarrow \frac{4}{\pi}.
\end{align*}

\begin{table}[h]
\centering
\caption{Log Family: $\text{PCF}(-kn^2,\,(k\!+\!1)n\!+\!k) = 1/\ln(k/(k\!-\!1))$}
\begin{tabular}{cccc}
\hline
$k$ & $a(n)$ & $b(n)$ & Value \\
\hline
2 & $-2n^2$ & $3n+2$ & $1/\ln 2$ \\
3 & $-3n^2$ & $4n+3$ & $1/\ln(3/2)$ \\
4 & $-4n^2$ & $5n+4$ & $1/\ln(4/3)$ \\
5 & $-5n^2$ & $6n+5$ & $1/\ln(5/4)$ \\
$k$ & $-kn^2$ & $(k\!+\!1)n\!+\!k$ & $1/\ln(k/(k\!-\!1))$ \\
\hline
\end{tabular}
\end{table}

\begin{table}[h]
\centering
\caption{Pi Family: $\text{PCF}(-n(2n\!-\!c),\,3n\!+\!1)$ for odd $c = 2m+1$}
\begin{tabular}{ccccc}
\hline
$m$ & $c$ & $a(n)$ & Value & $\text{val}\times\pi$ \\
\hline
0 & 1 & $-n(2n-1)$ & $2/\pi$ & $2$ \\
1 & 3 & $-n(2n-3)$ & $4/\pi$ & $4$ \\
2 & 5 & $-n(2n-5)$ & $16/(3\pi)$ & $16/3$ \\
3 & 7 & $-n(2n-7)$ & $32/(5\pi)$ & $32/5$ \\
4 & 9 & $-n(2n-9)$ & $256/(35\pi)$ & $256/35$ \\
$m$ & $2m\!+\!1$ & $-n(2n\!-\!(2m\!+\!1))$ & $\frac{2^{2m+1}}{\pi\binom{2m}{m}}$ & $\frac{2^{2m+1}}{\binom{2m}{m}}$ \\
\hline
\end{tabular}
\end{table}

\section{Discovery Methodology}

Both families were discovered via systematic \emph{modular signature sweeps}:
fixing the denominator polynomial $b(n) = kn + c$ for small $k$ and $c$,
then exhaustively searching quadratic numerator polynomials $a(n)$ using
a Meet-In-The-Middle hash-table algorithm at low precision, followed by
high-precision verification of hits.

The key insight was that the $k = 3$ family ($b(n) = 3n + c$) produced
$\pi$-related and $\ln$-related constants, in contrast to the classical
$k = 2$ family ($b(n) = 2n + 1$) associated with Brouncker's formula $4/\pi$.

\section{Numerical Certification}

\begin{table}[h]
\centering
\caption{Arb ball arithmetic certification (depths 2000--4000, 10000-bit precision)}
\begin{tabular}{lccc}
\hline
PCF & Target & Certified digits & Bracket width \\
\hline
$a = -2n^2,\, b = 3n+2$      & $1/\ln 2$     & $1208$ & $< 10^{-1208}$ \\
$a = -3n^2,\, b = 4n+3$      & $1/\ln(3/2)$  & $1817$ & $< 10^{-1817}$ \\
$a = -n(2n-1),\, b = 3n+1$   & $2/\pi$       & $1207$ & $< 10^{-1207}$ \\
$a = -n(2n-3),\, b = 3n+1$   & $4/\pi$       & $1217$ & $< 10^{-1217}$ \\
\hline
\end{tabular}
\end{table}

\noindent All four cases are certified by computing consecutive convergents
$C_N$ and $C_{N+1}$ using Arb interval arithmetic, with the true limit
rigorously contained in $[\min(C_N, C_{N+1}),\, \max(C_N, C_{N+1})]$.

\section{Open Questions}

\begin{enumerate}
\item \textbf{Proof of the Log Family.}
The hypergeometric connection via $\ln(1+x) = x \cdot {}_2F_1(1,1;2;-x)$
suggests that an equivalence transform of the Gauss continued fraction
for ${}_2F_1(1,1;2;z)$ should yield our PCF. Establishing this transformation
rigorously would convert Conjecture~\ref{conj:log} into a theorem.

\item \textbf{Proof of the Pi Family.}
The central-binomial closed form $2^{2m+1}/(\pi\binom{2m}{m})$
and its Gamma-function reformulation suggest a connection to the
Wallis product or to ${}_2F_1$ evaluations at specific arguments
involving cube roots of unity (given the $3n+1$ step).

\item \textbf{Non-integer parameters.}
The log family was verified for non-integer $k$ (including $k = 1.5, 2.5, \ldots$)
to $90$--$120$ digits. Does the pi family similarly extend to real $m$?

\item \textbf{Irrationality measures.}
Do these PCFs yield new irrationality measures for $\ln 2$, $\ln(3/2)$,
or $1/\pi$, following the approach of Ap\'{e}ry's theorem?
\end{enumerate}

\begin{thebibliography}{9}
\bibitem{raayoni2021}
G.\,Raayoni \textit{et al.},
``Generating conjectures on fundamental constants with the Ramanujan Machine,''
\textit{Nature} \textbf{590} (2021), 67--73.

\bibitem{elimelech2024}
D.\,Elimelech \textit{et al.},
``Algorithm-assisted discovery of an intrinsic order among mathematical constants,''
\textit{PNAS} \textbf{121}(25) (2024).

\bibitem{cuyt2008}
A.\,Cuyt, V.\,B.\,Petersen, B.\,Verdonk, H.\,Waadeland, and W.\,B.\,Jones,
\textit{Handbook of Continued Fractions for Special Functions},
Springer, 2008.

\bibitem{brouncker1656}
W.\,Brouncker (1656), continued fraction for $4/\pi$.

\bibitem{wallis1656}
J.\,Wallis, \textit{Arithmetica Infinitorum}, 1656.
\end{thebibliography}

\end{document}
